\documentclass[10pt]{article}

\usepackage[utf8]{inputenc}

\usepackage[T1]{fontenc}

\usepackage{amsmath}

\usepackage{amsfonts}

\usepackage{amssymb}

\usepackage[version=4]{mhchem}

\usepackage{stmaryrd}

\usepackage{bbold}



\begin{document}

рость света в вакууме ( $m_{0} c^{2}-$ энергия покоя точки). При малых скоростях $(v \ll c)$ последнее соотношение переходит в обычную формулу



\[

m_{0} v^{2} / 2

\]



C. M. Tapr.



См. также Энергия.



КИНЕТИЧЕСКИЙ МОМЕНТ - см. Момент количества движения.



КИНЕТИЧЕСКИЙ ПОТЕНЦИАЛ - см. Лагранжа фУнкЦИЯ. КИНЕТИЧЕСКОЕ УРАВНЕНИЕ - эволюцИонное уравнение для одночастичной функции распределения $F_{1}(t, x)$ $\left(x \equiv(q, p) \in \mathbb{R}^{v} \times \mathbb{R}^{v}-\right.$ координата и импульс частицы) типа



\[

\frac{d}{d t} F_{1}(t, x)=A\left(x ; F_{1}(t, x)\right)

\]



где выражение $A\left(x ; F_{1}(t, x)\right)$ в любой момент времени $t$ полностью определяется функцией $F_{1}(t, x)$ для этого же момента. Для квантовых систем К. у. определяется аналогичным образом для одночастичной редуцированной плотности матрицы (см. [1]). Основные примеры К. у.: Больцмана уравнение, Власова уравнение, Фоккера - Планка уравнение, Ландау уравнение, Больцмана - Энскога уравнение. Родственные им уравнения рассмотрены в [2].



Проблема обоснования К. у. связана с выводом К. у. из уравнений эволюции состояния статистич. систем - Боголюбова уравнений. Метод Н. Н. Боголюбова вывода К. у. состоит в том, что на основе асимптотич. разложений по малому параметру строится частное решение уравнений Боголюбова, к-рое зависит от времени через одночастичную функцию распределения. Если начальные функции распределения удовлетворяют принципу ослабления корреляций (см. Кластерное свойство), то одночастичные функции распределения определяются из К. у. Интерпретируется К.у. как уравнение эволюции состояния системы на определенном (кинетич.) этапе процесса релаксации к равновесному состоянию. Существуют и другие методы вывода K. у. (см. [2], [3]).



Так как состояние системы частиц описывается последовательностью $n$-частичных функший распределения, то, следовательно, К. у. описывает эволюцию состояния в определенной предельной ситуации. Поэтому в математич. работах последнего времени проблема обоснования К.у. формулируется как построение соответствующего предела для решений уравнений Боголюбова. Предельный переход совершается по параметрам, входяшим в гамильтониан системы, таким образом, что выполнение условия факторизации для предельных функций распределения в начальный момент времени влечет за собой выполнение аналогичного условия в любой момент времени (в системе отсутствуют корреляции). Следует подчеркнуть, что в кинетич. пределе описываются системы бесконечного числа частиц.



Предел решений уравнений Боголюбова удовлетворяет цепочке предельных уравнений Боголюбова. (Для стационарных решений уравнений Боголюбова см. Равновесных состояний кинетический предел.) Если в начальном состоянии отсутствуют корреляции, то состояние предельной системы в любой момент времени описывается одночастичной функцией распределения, к-рая удовлетворяет К.у. Таким образом, приходят к уравнению Больцмана в Больцмана - Грэда пределе (см. [3]), к уравнению Власова в среднего поля пределе и т. П. (см. [4]).



В настоящее время строгие результаты, связанные с выводом К. у., получены в пределе среднего поля для гладкого потенциала и интегрируемых начальных данных, в пределе Больцмана - Грэда для системы упругих шаров.



Вопросы сушествования и единственности решений К. у. весьма нетривиальны, известны случаи точных решений (см. [4]). К.у. находят широкое применение в различных разделах физики.



Лит.: [1] Б о го л ю бо в Н. Н., Избранные труды, т. 2, Киев, 1970; [2] Л и бо в Р., Введение в теорию кинетических уравнений, пер с англ., М., 1974; [3] Грэ д Г., в кн.: Термодинамика газов, пер. с англ. и нем., М., 1970, с. 5-168; [4] Итоги науки и техники. Современные проблемы математики. Фундаментальные направления, т. 2, M., 1985, c. 233-307.



В. И. Герасименко. KИНЕТИЧЕСКОЕ УРАВНЕНИЕ Б ОльцМана - см. Больимана кинетическое уравнение.



КИНК - термин, впервые возникший в теории синус-Гордона уравнения



\[

u_{t t}-u_{x x}+\sin u=0 \text {. }

\]



К инк $u_{+}(x, t)$ и антикинк $u_{-}(x, t)$ - это односолитонные решения уравнения синус-Гордона:



\[

u_{ \pm}(x, t)=4 \operatorname{arctg} \exp \left[ \pm\left(1-v^{2}\right)^{-1 / 2}\left(x-v t-x_{0}\right)\right] .

\]



Убывающие (по $\bmod 2 \pi$ ) решения уравнения синус-Гордона имеют важную целочисленную характеристику - топологич. заряд



\[

Q=\frac{1}{2 \pi} \int_{-\infty}^{\infty} \frac{\partial u(x, t)}{\partial x} d x .

\]



Если $Q=1$, то $u(x, t)-$ солитон (кинк, флюксон), если $Q=-1$, то $u(x, t)$ - антисолитон (антикинк, антифлюксон) (см. [1], [2]). Позднее, для других нелинейных эволюционных уравнений термин «К.» («антикинк») стал употребляться для одиночных солитонов, имеющих форму «перегиба». Напр., К. (антикинк) уравнения Миуры - Кортевега - де Фриса



\[

u_{t}-6 u^{2} u_{x}+u_{x x x}=0

\]



имеет вид:



\[

u(x, t)= \pm a \operatorname{th}\left[a\left(x+2 a^{2} t-x_{0}\right)\right] .

\]



См. также Беклунда преобразование.



Лит.: [1] Солитоны в действии, пер. с англ., М., 1981; [2] Т а х таджя н Л.А., Фадде ев Л. Д., Гамильтонов подход в теории солитонов, М., 1986.



В. П. Котляров.



KИРАЛЬНАЯ СИММЕТРИЯ в кван товой те о р и и п о л я - симметрия уравнений движения, которая комбинируется из двух различных симметрий: симметрии взаимодействия адронов относительно обычных преобразований в изотопическом пространстве без изменения внутренней четности и той же симметрии, но с изменением внутренней четности. Таким образом, преобразования К. с., кроме перемешивания состояний частиц с различными электрич. зарядами, перемешивают и состояния с разной внутренней четностью. К. с. является глобальной, то есть не зависяшей от точек пространства-времени. Такая инвариантность в случае частиц ненулевой массы не может быть связана ни с каким законом сохранения для фиксированной системы частиц, а определяет лишь форму их взаимодействия, напр. форму взаимодействия нуклонов с псевдоскалярными пионами, испускание каждого из к-рых изменяет четность системы. В этом смысле К. с. является динамической симметрией. К. с.- один из примеров симметрии, приводящей к существенно нелинейной квантовой теории поля.



Инвариантность относительно вращений в изотопич. пространстве без изменения четности связана с законом сохранения векторных токов $(V)$, а с изменением четности - с законом сохранения аксиальных токов $(A)$ (см. Ток в квантовой теории поля). Сохранение векторного тока можно связать с сохранением полного электрич. заряда системы взаимодействующих частиц. В случае безмассовых спинорных (со спином 1/2) частиц, напр. нейтрино, сохранение аксиального тока можно связать с определенным законом сохранения - законом сохранения спиральности. Действительно, в случае безмассового спинорного поля, распространяющегося со скоростью света, спин квантов поля направлен либо против движения, либо в сторону движения. Соответственно различают левую и правую спиральность; первому случаю соответствует комбинация $V-A$ токов частиц, второму - комбинация $V+A$, и эти комбинации должны сохраняться в отсутствие взаимодействия нейтрино с другими частицами. Однако если спинорная частица имеет ненулевую массу покоя, то ее спин не обязательно должен быть ориентирован по оси движения. Но во взаимодействиях с другими частицами это качество спиральности опять проявляется. Так, в слабом взаимодействии участвуют только лептоны с левыми спиральностями, а в





\end{document}